\documentclass[11pt]{article}

% Shared notes settings for Elements of AIML lectures (UPES).
% Keep this file free of \documentclass to allow reuse via \input.

\usepackage[utf8]{inputenc}
\usepackage[T1]{fontenc}
\usepackage[a4paper,margin=1in]{geometry}
\usepackage{hyperref}
\usepackage{xurl}
\usepackage{booktabs}
\usepackage{amsmath}
\usepackage{amssymb}
\usepackage{listings}
\usepackage{xcolor}
\usepackage{enumitem}
\usepackage{parskip}

% Code listings (general-purpose).
\lstset{
  basicstyle=\ttfamily\small,
  keywordstyle=\color{blue},
  commentstyle=\color{gray},
  stringstyle=\color{teal},
  showstringspaces=false,
  columns=fullflexible,
  frame=single,
  framerule=0.2pt,
  breaklines=true
}

\setlist[itemize]{topsep=0.3em,itemsep=0.2em}
\setlist[enumerate]{topsep=0.3em,itemsep=0.2em}

% Simple per-section source line.
\newcommand{\notesources}[1]{\par\noindent{\small\color{gray}\textit{Sources:} \sloppy #1\par}}

% Unit-wise source strings for this subject (used by slides/notes).
% Path is relative to the lecture `latex/` folder (the compilation CWD).
\IfFileExists{../../../_shared/aiml-sources.tex}{%
  % Source strings for Elements of AIML (used by slides + notes).

\newcommand{\AIMLRepoURL}{https://github.com/tali7c/Elements-of-AIML}

\newcommand{\AIMLLabSources}{%
Provost and Fawcett, \textit{Data Science for Business}.;%
McKinney, \textit{Python for Data Analysis}.;%
WEKA Documentation (Waikato Environment for Knowledge Analysis), accessed 2026-02-28.;%
Matplotlib and Seaborn documentation, accessed 2026-02-28.%
}

\newcommand{\AIMLUnitISources}{%
Rich and Knight, \textit{Artificial Intelligence}, McGraw Hill, 2017.;%
Russell and Norvig, \textit{Artificial Intelligence: A Modern Approach} (4e), Ch. 1--2 (Introduction, Intelligent Agents).;%
Poole and Mackworth, \textit{Artificial Intelligence: Foundations of Computational Agents} (2e), selected sections.%
}

\newcommand{\AIMLUnitIISources}{%
Russell and Norvig, \textit{Artificial Intelligence: A Modern Approach} (4e), Ch. 7--9 (Logical Agents, First-Order Logic, Inference).;%
Huth and Ryan, \textit{Logic in Computer Science} (2e), selected sections on propositional and first-order logic.;%
Genesereth and Nilsson, \textit{Logical Foundations of Artificial Intelligence}, selected chapters.%
}

\newcommand{\AIMLUnitIIISources}{%
Mueller and Massaron, \textit{Machine Learning for Dummies}, For Dummies, 2016.;%
Mitchell, \textit{Machine Learning}, selected chapters on learning basics and evaluation.;%
Scikit-learn documentation, accessed 2026-02-28.%
}

\newcommand{\AIMLUnitIVSources}{%
Mueller and Massaron, \textit{Machine Learning for Dummies}, For Dummies, 2016.;%
James, Witten, Hastie, and Tibshirani, \textit{An Introduction to Statistical Learning} (2e), selected chapters.;%
Scikit-learn documentation, accessed 2026-02-28.%
}

\newcommand{\AIMLUnitVSources}{%
NIST, \textit{AI Risk Management Framework (AI RMF 1.0)}, 2023.;%
Barocas, Hardt, and Narayanan, \textit{Fairness and Machine Learning} (online book), selected sections.;%
Knaflic, \textit{Storytelling with Data}, selected chapters on communicating results.%
}
%
}{%
  % Faculty copies live under `private/` so their compile CWD is different.
  \IfFileExists{../../../../public/_shared/aiml-sources.tex}{%
    % Source strings for Elements of AIML (used by slides + notes).

\newcommand{\AIMLRepoURL}{https://github.com/tali7c/Elements-of-AIML}

\newcommand{\AIMLLabSources}{%
Provost and Fawcett, \textit{Data Science for Business}.;%
McKinney, \textit{Python for Data Analysis}.;%
WEKA Documentation (Waikato Environment for Knowledge Analysis), accessed 2026-02-28.;%
Matplotlib and Seaborn documentation, accessed 2026-02-28.%
}

\newcommand{\AIMLUnitISources}{%
Rich and Knight, \textit{Artificial Intelligence}, McGraw Hill, 2017.;%
Russell and Norvig, \textit{Artificial Intelligence: A Modern Approach} (4e), Ch. 1--2 (Introduction, Intelligent Agents).;%
Poole and Mackworth, \textit{Artificial Intelligence: Foundations of Computational Agents} (2e), selected sections.%
}

\newcommand{\AIMLUnitIISources}{%
Russell and Norvig, \textit{Artificial Intelligence: A Modern Approach} (4e), Ch. 7--9 (Logical Agents, First-Order Logic, Inference).;%
Huth and Ryan, \textit{Logic in Computer Science} (2e), selected sections on propositional and first-order logic.;%
Genesereth and Nilsson, \textit{Logical Foundations of Artificial Intelligence}, selected chapters.%
}

\newcommand{\AIMLUnitIIISources}{%
Mueller and Massaron, \textit{Machine Learning for Dummies}, For Dummies, 2016.;%
Mitchell, \textit{Machine Learning}, selected chapters on learning basics and evaluation.;%
Scikit-learn documentation, accessed 2026-02-28.%
}

\newcommand{\AIMLUnitIVSources}{%
Mueller and Massaron, \textit{Machine Learning for Dummies}, For Dummies, 2016.;%
James, Witten, Hastie, and Tibshirani, \textit{An Introduction to Statistical Learning} (2e), selected chapters.;%
Scikit-learn documentation, accessed 2026-02-28.%
}

\newcommand{\AIMLUnitVSources}{%
NIST, \textit{AI Risk Management Framework (AI RMF 1.0)}, 2023.;%
Barocas, Hardt, and Narayanan, \textit{Fairness and Machine Learning} (online book), selected sections.;%
Knaflic, \textit{Storytelling with Data}, selected chapters on communicating results.%
}
%
  }{%
    % Source strings for Elements of AIML (used by slides + notes).

\newcommand{\AIMLRepoURL}{https://github.com/tali7c/Elements-of-AIML}

\newcommand{\AIMLLabSources}{%
Provost and Fawcett, \textit{Data Science for Business}.;%
McKinney, \textit{Python for Data Analysis}.;%
WEKA Documentation (Waikato Environment for Knowledge Analysis), accessed 2026-02-28.;%
Matplotlib and Seaborn documentation, accessed 2026-02-28.%
}

\newcommand{\AIMLUnitISources}{%
Rich and Knight, \textit{Artificial Intelligence}, McGraw Hill, 2017.;%
Russell and Norvig, \textit{Artificial Intelligence: A Modern Approach} (4e), Ch. 1--2 (Introduction, Intelligent Agents).;%
Poole and Mackworth, \textit{Artificial Intelligence: Foundations of Computational Agents} (2e), selected sections.%
}

\newcommand{\AIMLUnitIISources}{%
Russell and Norvig, \textit{Artificial Intelligence: A Modern Approach} (4e), Ch. 7--9 (Logical Agents, First-Order Logic, Inference).;%
Huth and Ryan, \textit{Logic in Computer Science} (2e), selected sections on propositional and first-order logic.;%
Genesereth and Nilsson, \textit{Logical Foundations of Artificial Intelligence}, selected chapters.%
}

\newcommand{\AIMLUnitIIISources}{%
Mueller and Massaron, \textit{Machine Learning for Dummies}, For Dummies, 2016.;%
Mitchell, \textit{Machine Learning}, selected chapters on learning basics and evaluation.;%
Scikit-learn documentation, accessed 2026-02-28.%
}

\newcommand{\AIMLUnitIVSources}{%
Mueller and Massaron, \textit{Machine Learning for Dummies}, For Dummies, 2016.;%
James, Witten, Hastie, and Tibshirani, \textit{An Introduction to Statistical Learning} (2e), selected chapters.;%
Scikit-learn documentation, accessed 2026-02-28.%
}

\newcommand{\AIMLUnitVSources}{%
NIST, \textit{AI Risk Management Framework (AI RMF 1.0)}, 2023.;%
Barocas, Hardt, and Narayanan, \textit{Fairness and Machine Learning} (online book), selected sections.;%
Knaflic, \textit{Storytelling with Data}, selected chapters on communicating results.%
}
%
  }%
}


\title{Elements of AIML\\Unit 01 -- Lecture 03 Notes\\AI Techniques, Limitations \& Impact}
\author{Tofik Ali}
\date{\today}

\begin{document}
\maketitle
\tableofcontents

\section{Lecture Overview}
This lecture consolidates Unit I by revisiting major AI technique families, discussing real-world limitations, and exploring how AI impacts society and application domains.
\notesources{\AIMLUnitISources}

\section{How to use these notes (self-study)}
This lecture is most useful if you treat it like a \textbf{real deployment review}.
For each example system, ask: (1) what is the task, (2) what data is used, (3) what can go wrong, and (4) what should we monitor after release.
Try to answer the practice questions without looking at the solutions first.

\section{Learning Outcomes}
\begin{itemize}[leftmargin=*]
  \item Recall core AI technique families and typical use cases.
  \item Explain major AI limitations in practice (data, drift, robustness, bias, explainability).
  \item Discuss AI's impact (benefits and risks) and identify application domains.
\end{itemize}

\section{Keywords}
\begin{itemize}[leftmargin=*]
  \item \textbf{Distribution shift:} when data in deployment differs from training distribution.
  \item \textbf{Robustness:} stability of performance under noise/perturbations.
  \item \textbf{Bias:} systematic error that disadvantages some groups or situations.
  \item \textbf{Explainability:} ability to provide understandable reasons for outputs.
  \item \textbf{Monitoring:} tracking model/data behavior in production for drift/failure.
\end{itemize}

\section{Abbreviations used in this lecture}
\begin{itemize}[leftmargin=*]
  \item \textbf{OOD}: out-of-distribution (inputs unlike training data).
  \item \textbf{PII}: personally identifiable information (privacy-sensitive data).
  \item \textbf{KPI}: key performance indicator (business/operational metric).
  \item \textbf{Drift}: changing data patterns over time (often a form of distribution shift).
\end{itemize}

\section{AI technique families (recap)}
\subsection{Search and optimization}
Many AI problems can be framed as searching a space of possible solutions:
\begin{itemize}[leftmargin=*]
  \item path planning, scheduling, game playing, constraint satisfaction.
\end{itemize}

\subsection{A quick mental model: problem $\rightarrow$ technique}
\begin{itemize}[leftmargin=*]
  \item If you have \textbf{rules/constraints}: use logic/constraints/search.
  \item If you have \textbf{labeled data}: use supervised learning.
  \item If you have \textbf{unlabeled data}: use clustering/dimensionality reduction.
  \item If you have \textbf{interaction + reward}: use reinforcement learning.
\end{itemize}

\subsection{Knowledge and logic}
Logic-based AI stores explicit knowledge (facts and rules) and uses inference. This is the focus of Unit II.

\subsection{Machine learning}
ML learns patterns from data, which is the focus of Units III--IV. ML is powerful when rules are hard to write but examples are available.

\section{Limitations and failure modes}
\subsection{Data limitations}
Practical AI/ML systems rely on data that may be:
\begin{itemize}[leftmargin=*]
  \item incomplete (missing values), noisy, biased, or incorrectly labeled,
  \item not representative of real deployment.
\end{itemize}

\subsection{Data quality checklist (practical)}
\begin{itemize}[leftmargin=*]
  \item Are there missing values? Are they random or systematic?
  \item Are labels reliable? Who labeled the data and with what guideline?
  \item Are there duplicates/near-duplicates that can leak into the test set?
  \item Do some groups/users/regions appear rarely (representation gaps)?
\end{itemize}

\subsection{Distribution shift}
Even if a model performs well on a test set, it may fail later when the world changes:
\begin{itemize}[leftmargin=*]
  \item new user behavior, new policies, new markets, or adversarial changes.
\end{itemize}

\subsection{Types of distribution shift (intuition)}
It helps to distinguish:
\begin{itemize}[leftmargin=*]
  \item \textbf{Covariate shift:} input distribution changes ($P(X)$ changes), e.g., different camera conditions.
  \item \textbf{Label shift:} class proportions change ($P(Y)$ changes), e.g., fraud rate increases during a festival.
  \item \textbf{Concept drift:} relationship changes ($P(Y|X)$ changes), e.g., attackers change strategies so the same features mean different things.
\end{itemize}

\subsection{Robustness and adversarial behavior}
Small changes in input can sometimes cause large output changes. In security settings, attackers may deliberately manipulate inputs to fool a system.

\subsection{Robustness vs security}
\begin{itemize}[leftmargin=*]
  \item \textbf{Robustness} is about natural variation/noise: typos, lighting changes, sensor noise.
  \item \textbf{Security} is about intentional manipulation: adversarial inputs, fraudsters, prompt injection/misuse.
\end{itemize}
Both require test cases beyond the standard i.i.d. (independent and identically distributed) assumption.

\subsection{Explainability and accountability}
In high-stakes decisions (loans, healthcare, admissions), stakeholders often require:
\begin{itemize}[leftmargin=*]
  \item reasons for decisions,
  \item ability to contest errors,
  \item documented limitations.
\end{itemize}

\subsection{Explainability vs interpretability (short)}
\begin{itemize}[leftmargin=*]
  \item \textbf{Interpretability:} the model is inherently understandable (e.g., small decision tree, linear model).
  \item \textbf{Explainability:} tools that provide reasons/attributions for a model's outputs (even if model is complex).
\end{itemize}
In practice, we often use a combination: interpretable baselines + explanation tools + documentation.

\subsection{Safety, privacy, and misuse}
AI can cause harm if used without safeguards:
\begin{itemize}[leftmargin=*]
  \item privacy leakage, unsafe automation, surveillance misuse.
\end{itemize}

\subsection{Privacy: data and model}
Privacy risks are not only about datasets:
\begin{itemize}[leftmargin=*]
  \item datasets may contain \textbf{PII} (names, IDs, medical information),
  \item models can leak information about training data (membership inference/memorization),
  \item logs and analytics can unintentionally expose sensitive data.
\end{itemize}

\section{Deployment checklist (practical)}
\begin{enumerate}[leftmargin=*]
  \item Define the goal and the real cost of errors (who is harmed by mistakes?).
  \item Validate that training data matches deployment and identify shift risks.
  \item Evaluate fairness across relevant groups and document limitations.
  \item Plan monitoring: drift detection, retraining cadence, and rollback strategy.
\end{enumerate}

\subsection{What to monitor (concrete signals)}
\begin{itemize}[leftmargin=*]
  \item \textbf{Data drift:} missing rate, min/max/mean changes, outlier frequency, category frequency changes.
  \item \textbf{Prediction drift:} score distributions, confidence, class proportions predicted by the model.
  \item \textbf{Performance:} accuracy/precision/recall when labels arrive (possibly delayed).
  \item \textbf{Operational:} latency, cost, error logs, user feedback.
\end{itemize}

\section{Impact and application domains}
AI has broad impact:
\begin{itemize}[leftmargin=*]
  \item \textbf{Benefits:} improved decision support, efficiency, personalization, new services.
  \item \textbf{Risks:} bias amplification, privacy loss, over-reliance, job changes.
\end{itemize}
Application domains include finance, healthcare, security, education, retail, logistics, energy, and more.

\subsection{Responsible practice mindset}
Responsible use is not ``one extra checkbox''; it is a set of ongoing practices:
\begin{itemize}[leftmargin=*]
  \item define stakeholders and harms,
  \item choose metrics that reflect real costs,
  \item test across groups and scenarios,
  \item document assumptions and limitations,
  \item monitor and update after deployment.
\end{itemize}

\section{Mini case study: credit scoring}
Consider an ML model for loan approval.
\subsection{What can go wrong?}
\begin{itemize}[leftmargin=*]
  \item \textbf{Metric trap:} optimizing accuracy while ignoring unequal costs of errors.
  \item \textbf{Bias:} training data encodes historical unfairness.
  \item \textbf{Proxy features:} some variables correlate with protected attributes.
  \item \textbf{Drift:} economy changes, making the model outdated.
\end{itemize}
\subsection{What to do?}
\begin{itemize}[leftmargin=*]
  \item use cost-aware metrics, validate fairness, add monitoring, and document decisions.
\end{itemize}

\section{Common pitfalls (exam-focused)}
\begin{itemize}[leftmargin=*]
  \item \textbf{Wrong metric:} optimizing accuracy when the real cost is dominated by false negatives or false positives.
  \item \textbf{Assuming ``train = real world'':} ignoring distribution shift (OOD inputs, drift, seasonality).
  \item \textbf{No monitoring plan:} treating deployment as the end, not the beginning (no drift/performance tracking).
  \item \textbf{Robustness vs security confusion:} robustness is about noise/shift; security includes adversarial behavior and misuse.
\end{itemize}

\section{Practice (with answers)}

\subsection*{P1. Identify shift type.}
Give the shift type for each scenario:
\begin{itemize}[leftmargin=*]
  \item New camera sensor changes brightness distribution for a vision model.
  \item Fraud becomes more common during a holiday season.
  \item Attackers change wording to bypass a spam classifier.
\end{itemize}
\textbf{Answer:} camera brightness change $\rightarrow$ covariate shift; fraud rate change $\rightarrow$ label shift; attackers changing strategy $\rightarrow$ concept drift.

\subsection*{P2. Monitoring design.}
List two signals you would monitor for a deployed spam classifier.
\textbf{Answer:} predicted spam rate over time; distribution of top features (e.g., URL count); user ``not spam'' feedback rate; latency and error logs.

\section{Exit questions (with answers)}

\subsection*{Q1. List four reasons why AI systems fail in deployment.}
\textbf{Answer:} bad/non-representative data, distribution shift, robustness issues, leakage, wrong metric for the goal, no monitoring, bias/fairness issues.

\subsection*{Q2. Give two examples of distribution shift.}
\textbf{Answer:} (1) user behavior changes after a new policy; (2) fraud patterns evolve; (3) sensor changes or new device types.

\subsection*{Q3. In a credit scoring system, why is explainability important?}
\textbf{Answer:} decisions affect people; stakeholders need justification and accountability; regulations and trust require understandable reasons and error correction.

\subsection*{Q4. Name five AI domains and one task in each.}
\textbf{Answer:} healthcare (risk prediction), finance (fraud detection), security (anomaly detection), education (recommendation), logistics (route optimization), retail (demand forecasting), energy (load forecasting), public services (resource allocation).

\section{Summary}
\begin{itemize}[leftmargin=*]
  \item AI combines search, logic, and learning.
  \item Real-world success requires careful evaluation and responsible deployment.
\end{itemize}
\notesources{\AIMLUnitISources}

\end{document}
